\RequirePackage[T1]{fontenc}
\documentclass[conference,10pt]{IEEEtran}
\usepackage{amsmath,amssymb,graphicx,url}
\title{Current-STA Transactional Repair for Timing--Resource Tradeoffs in 3D FPGA Routing}
\author{\IEEEauthorblockN{Anonymous development manuscript}}
\begin{document}
\maketitle
\begin{abstract}
Legal FPGA routes can occupy different points in a timing--resource design space. In a fixed 3D routing resource graph, sharing a vertical resource may save occupancy but restrict a critical connection or create capacity competition. We study a current-STA repair layer that selects complete connection groups, performs native searches on private route-tree copies, and commits only after physical, capacity and native timing checks. Retained-crossing domains encourage sharing; a failed guided search can release only its network's domain while preserving successful restorations elsewhere. Development results include three automatic timing repairs of a strong native BNN route, reducing critical-path delay from13.607436 to13.228146 ns with forty additional planar wire units. Subsequent resource refinement exposes distinct low-wire and low-vertical-occupancy alternatives at the same delay. Complete resource-first timing repair reaches13.228146 ns with43 SB and5060 CB, versus82 SB and5062 CB for its native control, at a cost of1762 more wire units. An independent five-connection feedback intervention improves ten setup components without regressions. On eltwise, a uniform K8 grouping cap reduces actual SB occupancy to24 and CB occupancy to965 at unchanged delay; an R1-based SPMV repair remains dominated by strong A0. A fixed-seed R1 resource trajectory reaches33 SB at unchanged13.421610 ns, spending1721 more wire units. Matched native controls distinguish selection, grouping and domain contributions. Prototype accounting exposes construction and validation overhead; complete cold algorithm timing and broader fixed-policy confirmation remain open.
\end{abstract}
\begin{IEEEkeywords}
3D FPGA, detailed routing, static timing analysis, transactional repair, multiobjective optimization
\end{IEEEkeywords}

\section{Problem and Motivation}
Three-dimensional FPGA fabrics couple directed reachability, scarce vertical channels and detailed timing. A different route can improve timing, save resources or exchange one benefit for another even when packing, placement and architecture remain fixed. A repair that does not reduce the final critical-path delay (CPD) may still be useful under an explicit resource or timing requirement. Conversely, a locally shorter connection is not evidence of a better complete route: another endpoint can become critical, or the replacement may prevent a different network from using a scarce channel.

VTR provides the detailed routing and timing infrastructure used here~\cite{vtr9paper}. The placement and architectural context follows the separation of design tradeoffs emphasized by Escaping Flatland~\cite{flatland}; our comparisons hold each physical input fixed. This work concerns post-legal routing decisions. It neither changes the physical timing model nor infers silicon area or power from reduced occupancy on a fixed device.

Initial shared-crossing guide experiments motivated the change of intervention point. Candidate coverage, search adoption, retained branches and native convergence can disconnect a guide's proxy score from final timing. We therefore operate on the current routed state and native STA, and evaluate an actual replacement rather than assuming that a different plan is adopted. The repair search itself remains native. The contributions under investigation are group selection, dependency scope, retained-domain decisions, acceptance and scheduling.

Our evidence distinguishes three questions: does an actual legal route improve a joint quality coordinate; which mechanism caused the change under matched local controls; and what does a complete fixed configuration cost? The first two have development evidence. The last requires further controlled full-flow measurements, rather than substituting a repair core or a shared research workflow.

\section{Objective and State Identity}
An input identity fixes the directed resource graph, packed netlist, placement, endpoint mapping, SDC and clock domains. A legal route $R$ uses actual graph edges and respects node capacities and native routing constraints. We record
\begin{equation}
 q(R)=(D(R),W(R),S(R),C(R),T(R)),
\end{equation}
where $D$ is native final CPD, $W$ is planar wire in the auditor's native resource units, $S$ is occupied SB pairs, $C$ is occupied CB output sources, and $T$ is comparable complete algorithm time. Paired SB half-nodes are one vertical resource; multiple CB edges from one output are not independent outputs. Guide recommendations are not occupancy measurements.

A final legal candidate is retained when it is not dominated on the measured, comparable coordinates. Quality projections use the first four coordinates and preserve uncertainty from printed CPD. They do not establish five-dimensional dominance when complete time is unknown. Different convergence rules, repair opportunities and implementation versions remain explicit, even for a descriptive comparison on the same physical inputs.

Final retention is separate from a branch's transition rule. We retain a strict CPD mode, $D(R')<D(R)$, and a fixed-resource mode that reduces CPD within registered absolute resource limits. A resource-under-timing mode permits wire tradeoffs within a fixed envelope while requiring nonincreasing selected resource objectives and at least one strict reduction. The current resource probes select $(S,C)$ as those objectives, cap CPD at the root value, cap wire at the root value plus0.1\%, and prohibit increases beyond root vertical counts. These are design parameters, not general termination conditions.

An endpoint-vector control fixes all endpoint/launch/capture identities and compares complete native setup slack values. It accepts only nonworse CPD, no setup-component regression and at least one strict setup improvement, with no increased hold violation. Native numeric values are compared without an epsilon. Resource policies impose their stated CPD and hold conditions; they do not silently inherit the stricter endpoint-vector requirement. The new bounded vector variant additionally enforces the strict timing control's immutable wire/SB/CB limits. Anchors do not roll forward, and no tolerance can accumulate timing degradation.

\section{Mechanism: Current-State Whole Transactions}
\subsection{Select targets and dependency scopes}
Timing selection follows the current native winning-path structure, including tied fan-ins, rather than a permanently retained list of endpoints. Capacity queries distinguish a promising extension from one blocked by other networks. Owner scopes and grouped targets can release dependencies that a single-connection repair leaves frozen.

Resource selection scans the actual current route trees. A candidate removes all sinks below a selected crossing subtree, preserves paths to the other sinks, and requires both structural release potential and a retained crossing alternative. The screen uses current slack margin; it is not a proof of final resource savings. Ranking is deterministic: released crossings per removed sink, fewer sinks, larger current margin, then network and pin IDs. A K8 group contains at most eight distinct networks; fewer are used when fewer qualify. Removed connections are restored in descending current raw criticality with deterministic ties. Accepted states receive a fresh current selector.

\subsection{Native searches and retained-crossing domains}
For each group we copy the selected trees and prune exactly the registered sink sets. Occupancy is updated as every local replacement is installed, so later searches see earlier resource use. Native weighted search uses the full device search box and the existing timing/congestion objective. Capacity-infeasible results are not committed. Real switch delays and final STA remain unchanged.

The domain control initially allows cross-layer edges only through crossing sources retained after pruning. It preserves useful sharing but can become infeasible under current capacities. The unrestricted native control receives the same current source, sink set, tree scope, objective and search range. Paired development probes evaluate both variants and use the registered domain-first acceptance order; native-only trajectories instead reselect from their own accepted unrestricted states. Those trajectories must not borrow a guided successor after divergence.

Automatic failure release makes the domain a revisable decision. If guided native search fails to return a capacity-feasible branch, search costs are reset and only that network's domain is removed. Previously successful restorations and their occupancy stay in the private transaction. Other networks remain guided. Search failure is an operational trigger, not itself a mathematical infeasibility certificate. A separate capacity diagnosis can justify a stronger statement for a particular fixed scope.

\subsection{Complete checks, acceptance and restoration}
All selected connections must be restored before acceptance. Independent checks validate real edges, unchanged retained branches and terminals, actual global capacities and resource counts. Native FULL and fresh STA evaluate the actual replacement. Acceptance commits the entire group; failure or rejection retains the source state. Intermediate private states may spend resources or slack when a complete mechanism-backed combination is subsequently checked. Requiring every component to improve CPD would exclude some useful groups.

The initial prototype serializes the candidate and performs a separate native LOAD/FULL/STA. A subsequent independent implementation temporarily installs complete candidate copies in the same private worker, runs actual native FULL and a new timing-analysis object, and exports a dedicated candidate route and timing view. It then restores original trees, sink lookups, occupancy, route digest and output hooks, and independently reobserves the source. The worker's original-input timing is never relabeled as candidate timing. This integration has passed small-graph tests and the actual BNN resource transaction: dedicated candidate exports exactly equal a separate same-binary LOAD/FULL/STA, independent physical/scope checks pass, and source recovery is exact. A further same-binary paired export experiment also validates an actual free-path CPD-changing transaction against independent candidate LOAD/FULL/STA. Large CPD-changing forced-path transactions remain outside that validated scope.

Loaded post-legal contexts reconstruct occupancy and native STA but reset the original route run's accumulated PathFinder price history. Matched native/domain controls share this context. An eventual in-router integration must explicitly preserve or reset that history and evaluate the choice.

\section{Overall Development Results}
The implementation derives from VTR revision \texttt{0b912334fdd83bb754b302146ac0d1ec9edc9a0e}. Results below use fixed-input serial clustered routing on the evaluated two-layer, edge-encoded Hybrid-O graphs, ideal clock modeling and native FULL/STA. Flat routing, CHANZ resources and arbitrary multilayer instances are outside this implementation evidence. Each record binds the exact graph, packing, placement, constraints, binary, requests and acceptance rule. A0 is native routing; R1 is the older guided method. Strong BNN A0/R1 roots additionally allow eight native convergence attempts with a1.01 reconvergence threshold. BNN has no supplied SDC: native default constraints use one ideal ap\_clk domain, zero input/output delay and zero setup/hold constraints; exact constraint echoes are bound. Archived U and Tf predecessors remain separate input/version/convergence records, not interchangeable points in a new algorithm's aggregate score.

\begin{table}[t]
\centering\caption{BNN legal quality points. Fast A0 follows three automatically selected timing commits from pristine strong A0. Resource phases have their own explicitly recorded opportunities. This is a quality comparison, not a matched-compute algorithm ranking.}
\label{tab:bnn}
\small
\begin{tabular}{lrrrr}\hline
State & CPD(ns) & Wire & SB & CB\\\hline
Strong A0 &13.607436&2059567&94&5062\\
Strong R1 &13.421610&2068908&86&5072\\
A0 native K8 (end) &13.607436&2059591&82&5062\\
A0 guided K8 (end) &13.607436&2061353&43&5060\\
Fast A0/native K8 &13.228146&2059603&86&5062\\
Fast A0/guided K8 &13.228146&2059882&79&5061\\
Fast R1/guided K8 &13.228146&2069040&74&5071\\
Full native + timing &13.228146&2059627&82&5062\\
Full guided + timing &13.228146&2061389&43&5060\\
45SB vector feedback &13.599798&2061253&45&5060\\
R1 selective release &13.228146&2069104&72&5071\\
R1 fixed-seed K8 &13.421610&2070629&33&5071\\\hline
\end{tabular}
\end{table}
The strong-A0 timing trajectory accepts three current-state transactions and reduces CPD by2.787376\%, adding forty wire units with SB/CB unchanged. Its third step improves CPD but regresses a noncritical setup component, making the distinction between strict CPD and endpoint-vector policies observable. The same-source fixed-vertical policy independently selects the identical three states using exact completed-query/transaction reuse; its replay time is not evidence of speedup.

From fast A0, historical unrestricted K8 reaches86 SB at2059603 wire; guided K8 spends279 more wire for seven fewer SB and one fewer CB at the same CPD. Corresponding pristine-root two-accept resource prefixes preserve the slower13.607436 ns CPD and use forty fewer wire units. They are fixed historical opportunities, not converged native bounds. A separately sealed resource-first timing phase starts from79SB/5061CB, keeps these vertical caps and the original root wire envelope, and automatically reaches13.421610,13.413972 and13.228146 ns. The last repair expands a current critical-owner dependency to nine networks and twelve connections: one noncritical connection spends about1.114955 ns of route margin while a critical connection saves about0.185826 ns. Strict CPD accepts it; the endpoint-vector control rejects a setup regression. Across this phase eleven setup components improve and seventeen regress, with no increased hold violation; the fixed source near-critical band has one improvement and no regression. The final current-state portfolio has no further admissible proposal.

The resulting four quality coordinates and complete setup/hold endpoint vectors exactly equal the existing timing-first/guided-K8 endpoint, although the serialized route body differs. This is convergence to a known quality point, not a new frontier gain or a general proof that phase order commutes. All ten current winning-connection switch distances equal already-computed exact source-only optimistic references after graph/root/sink identity checks; no queries are repeated for this diagnosis. The reference relaxes occupancy, shared branches and native arithmetic, so it is not a global native-STA bound. These observations favor resource/runtime and other-scope interventions over merely increasing the same cone query budget. Such a candidate would be lost under a universal strict-CPD retention rule.

\begin{table}[t]
\centering\caption{Eltwise fixed-seed K8 controls from the same pristine A0 root, using uniform grouping and acceptance. Both branches use their own successors.}
\label{tab:eltwise}
\small
\begin{tabular}{lrrrr}\hline
State & CPD(ns) & Wire & SB & CB\\\hline
Strong A0 root &6.018101&217065&30&968\\
Native K8 &6.018101&217050&26&968\\
Paired fixed-seed K8 &6.018101&217183&24&965\\\hline
\end{tabular}
\end{table}
Eltwise guided K8 uses actual group sizes six then three. It saves one additional SB relative to the prior guided K4 endpoint, with identical CPD, wire and CB. No failure release triggers in this circuit, so the gain is attributed to grouping/domain decisions rather than fallback. The native-only branch accepts one move, then rejects three legal but resource-worse alternatives; these rejected candidates remain recorded. On SPMV, paired and native-first R1 repair produce the same legal route at7.045986461 ns,148605 wire,zero SB and967 CB, but remain quality-dominated by strong A0. Initial routability and a circuit-balanced timing improvement across the suite have not been established by these repair probes.

\section{Pareto Tradeoffs and Attribution}
\begin{figure*}[t]
\centering\includegraphics[width=.98\textwidth]{data/joint_quality_projection_10.pdf}
\caption{Archive388 contains267 BNN observations on one physical-input identity, including strong A0/R1/U/Tf, complete resource-first timing, independent45SB feedback and all legal candidates of the five completed fixed-seed resource roles. Duplicate quality tuples are collapsed to156 plotted coordinates. Distinct roots, convergence rules and opportunities remain in the accompanying data. These are not two-dimensional Pareto boundaries: other resource coordinates and complete time are hidden.}
\label{fig:joint}
\end{figure*}
Table~\ref{tab:bnn} and Fig.~\ref{fig:joint} retain external strong baselines. The fast A0 native-K8 point quality-dominates the earlier slower A0 guided-K4 point, which had the same86 SB and5062 CB but wire2059667. Slower native-first low-wire states and the CB5059 reassignment state retain distinct resource tradeoffs; the new timing prefix quality-dominates several CB5060 predecessors. Quality domination is not sufficient for a complete-runtime claim when comparable time is missing.

\subsection{Extending the native control changes the comparison}
The completed native-order resource cohort ends at13.607436067 ns for both own-current branches. The unrestricted control makes five accepts, using2059591 wire,82 SB and5062 CB. The paired domain/failure-release branch makes nine accepts, using2061353 wire,43 SB and5060 CB. Its additional1762 wire units buy39 fewer SB and two fewer CB. Both exhaust their registered current menus; historical two-accept points are prefixes, not native resource bounds. Native and paired workflows cost2061.28 and7293.64 s and create nine and nineteen native children respectively. The paired branch pays its additional searches; these imported-state measurements are not equal-compute or cold whole-configuration timing.

Continuing from each completed resource endpoint, three strict current-STA commits reach the identical13.228145690 ns CPD. The native branch uses2059627 wire/82 SB/5062 CB; the paired branch uses2061389/43/5060. The extra1762 wire units still buy39 SB and two CB, now at the faster timing point. Both are nondominated in the collected four-axis population and both timing menus terminate. Full-vector comparison has550 setup improvements and799 regressions, so equal CPD does not imply endpointwise dominance. The native imported timing phase costs3093.90 s; the paired recovery segment costs450.11 s, with original interrupted elapsed time unknown. Those durations are not comparable complete costs.

A separate historical composition starts from the earlier seven-accept45-SB/5060-CB prefix. Three strict repairs reach13.228145690 ns at2061301 wire/45 SB/5060 CB in3028.24 s of imported repair. This prefix is not the final43-SB resource input. Its full-vector branch accepts none of six proposals in2747.93 s: one target improves physically but leaves the endpoint vector unchanged because another tied path remains; larger combinations worsen donors.

Same-source rejection feedback removes actions whose donors increased delay, then jointly covers two tied-path targets with three donors. The actual five-connection transaction reaches13.599797732 ns, a7.638334 ps gain, at2061253 wire/45 SB/5060 CB. All ten changed setup components improve, none regress, and no hold violation increases; both unchanged strict and full-vector policies accept. The434.52 s intervention includes preparation, native coordination and candidate evaluation. Its lower wire use retains a distinct four-axis nondominated point despite slower CPD than the43-SB branch. This validates one feedback-selected combination, not a general controller or independent circuit sample.

\subsection{Automatic transfer exposes retained-tree order sensitivity}
With delay/RR-ID repair seed order fixed across inputs, the eltwise native/domain roles end at217050/217183 wire,26/24 SB and968/965 CB at equal6.018101129 ns; imported workflows cost123.75/126.10 s. BNN A0 ends at2059591/2061377 wire,82/43 SB and5062/5060 CB, both at13.607436067 ns, costing2076.31/7099.32 s. Compared with native-order predecessors, paired wire decreases by8 on eltwise but increases by24 on BNN A0, with CPD and vertical counts unchanged. Five of eight registered root/control roles are complete on two circuits; multiple roots and trajectory observations are not independent circuit samples.

The earlier217214-to217191-wire eltwise change comes from source traversal order. First-step physical trees/vectors and second requests/domains match. A same-binary intervention uses canonical source, native-serialized source and a native source with only6719 restored to canonical order; complete candidates/vectors reproduce old/new/old under FULL/freshSTA and recovery checks. Seed contents agree but traversal order differs. This does not prove equal weighted costs or a heap tie. Ordered bytes remain in search/cache identity.

\subsection{An explicit seed policy removes the observed ambiguity}
A new independent implementation orders eligible repair seeds either by RR ID or by native physical delay followed by RR ID; the default keeps native depth-first order. It preserves seed membership, native primary cost formulas, physical delays and capacity/domain checks. In six actual eltwise source-by-policy transactions, both explicit policies produce identical complete candidate trees, endpoint vectors and real sink-cost traces from the canonical and native-serialized sources. RR-ID order yields217214 wire; physical-delay/RR-ID order yields217183, with24 SB,965 CB and unchanged6.018101129 ns CPD. The latter saves eight wire units over the native-serialized result and thirty-one over the canonical result, passing the resource-under-timing rule. This local intervention and the separately completed pristine-root roles remain distinct evidence.

Instrumentation records finite sink labels before cost reset; all lookahead factors are zero. Three later cost rows differ, including pressure retries, so equal-cost ties alone do not explain the result. Twelve native and six tiny regression jobs validate defaults, seed contents, FULL/freshSTA, independent LOAD and restoration. The automatic A0/R1/Tf batch fixes the delay/RR-ID policy globally for native/domain controls; policies are not selected by circuit history.

The same globally fixed delay/RR-ID seed policy also completes an R1-rooted paired K8 resource trajectory. Eight commits reduce86 SB/5072 CB to33 SB/5071 CB while preserving13.421610490 ns; wire increases by1721 to2070629, within the fixed root envelope. Its37-SB intermediate and33-SB final point are distinct nondominated quality coordinates. At rounds4--7, matched unrestricted candidates from the paired branch's current states fail the resource rule while guided candidates pass; this establishes local domain contributions. The own-current R1-native trajectory remains checkpointed, so full-trajectory superiority is unproven. All eighteen legal candidates, including degraded and rejected alternatives, remain in the archive. The sealed imported workflow takes5046.61 s and nineteen native children; complete cold time is unknown.

\subsection{Capacity feedback changes the domain decision}
An all-eight-domain R1 transaction fails at network31767/pin1 after five successful restorations. A directed path exists when occupancy is ignored. Under the recorded transient occupancy, the retained domain has no capacity path while an unrestricted witness exists. Reconstructing the minimum occupancy immediately after pruning also leaves no domain-capacity path. Because subsequent restorations only add occupancy, changing order or search pressure cannot fix this particular frozen scope/domain. The diagnostic searches a relaxed superset of native re-expansion/closure choices; its negative result is a necessary infeasibility certificate for that scope, not for the entire FPGA.

Releasing only that network's domain yields a fully legal state at72 SB, compared with74 SB for the same-source unrestricted control, at the cost of64 wire units and unchanged CPD/CB. The automatic native release implementation obtains exactly the prior selective-domain route and full endpoint vectors without a manually supplied release network or restarting earlier restorations. Nine setup components improve and none regress, so endpoint-vector acceptance succeeds while strict CPD rejects. All23 components in a fixed source near-critical band stay unchanged. This step demonstrates resource and secondary-timing benefit, not later CPD improvement.

\subsection{Grouping contribution and residual timing potential}
The first native-only fast-A0 K8 group produces exactly the actual route and endpoint vectors of two K4 steps. The first K4 component alone is endpoint-flat and rejected; the complete K8 group improves one setup component with no regressions. This supports complete combination transactions without demanding that every intermediate component pass the final policy. It does not show that a pristine-root endpoint policy reaches that state.

The fast-A0 timing selector eventually has no positive winning-path action. We tested a broader diagnostic: remove all retained-prefix constraints and background occupancy, then query source-root distances for every one of its ten current winning connections. Seven changed queries are newly solved and three exact matches reused. All ten optimistic static-switch distances equal the current values, and every directed witness is independently checked. This limits the evidence for merely increasing the same timing-query budget or rerouting a shared prefix at this state. It is an exact reference for the verified zero-RC switch-delay relaxation, not a full native-arithmetic/STA global optimum or a limit on the overall research scope.

\section{Overhead, Limitations and Next Confirmation}
The automatic selective-release experiment takes287.624 s from its sealed imported source through final context checks. Its native repair core is6.693 s, nested inside a113.480 s worker request. Source materialization takes21.221 s, transaction/domain audits41.032/10.984 s, candidate FULL/STA and materialization65.870 s, final sealing14.856 s, and registration/context/miscellaneous work20.182 s. These stages cannot be added again to the nested native child times. The result motivates state and validation integration; it does not establish lower complete routing time.

Fused candidate observation passes ten native regression jobs and actual BNN resource/timing checks against separate same-binary LOAD/FULL/STA. Candidate and restored-source receipts remain distinct. Optional graph export preserves native search, ordered routes, full frontiers and restoration in ten tiny jobs and two BNN paired transactions. Full/compact requests take86.73/67.90 s for timing and90.21/70.66 s for resource, with reversed order. Compact mode omits2.21 GiB of unused graph text; timing-cone selection retains full output. These single component pairs exclude startup and validation-only LOADs and do not establish complete-flow speedup. Earlier recovery costs and decoder failures remain in the evidence archive.

The new45-SB timing composition further exposes integration cost: its first committed round has843.38 s of completed top-level child work, with a0.667 s native repair core nested inside transaction coordination and206.49 s of complete candidate analysis/materialization. These nested figures are not added together or presented as complete single-configuration time.

Current evidence has several practical limits. Timing improvements are concentrated in BNN development trajectories; resource transfer does not imply initial routing-success improvement. K1/K4/K8 and scheduling pilots are not one frozen configuration. Complete cold routing-plus-repair time remains unknown, and old U timing, search-core time or sums of historical stages cannot fill it. Confirmation must freeze circuit-independent selection and strategy parameters, start automatically from sealed Tf and strong native roots, and compare A0/R1/native repairs with matched domains, allowed tree scope, target selection and acceptance opportunities. Equal-source local controls provide causal evidence; they do not substitute for that whole-algorithm evaluation.

\subsection{An exact streaming decoder improvement}
A streaming reader validates every pin row while retaining only the required endpoint and metric data. On two357.20 MiB source/candidate traces, old/new times are11.48/5.10 s and10.80/5.31 s with reversed order; peak process memory falls from1.73 GiB to0.18--0.21 GiB. Full vectors, all52,495 setup components and endpoint/domain identities agree exactly. Native COMPLETE pin counts detect missing rows. Corruption checks and bound integration evidence remain in the supplement. These are component measurements, not routing speedups.

\subsection{Native counting targets measured validation overhead}
The completed extended native BNN imported-state workflow costs2061.28 s. Native worker requests account for572.64 s; candidate materialization/physical audit, transaction reconstruction and tree-equivalence checks account for543.80,330.43 and136.65 s. These exclusive categories explain why reducing graph text alone does not resolve prototype overhead. A new opt-in C++ census counts actual planar wire, globally occupied SB pairs and CB output sources, per-net occupancies and crossing edge kinds from actual route trees. In source/candidate/restored and ordinary-input observations, every count matches an independent physical audit. Counter on/off preserves ordered route bodies, complete native frontiers and seed/cost traces; an additional same-binary candidate LOAD independently passes physical and timing checks.

These counts carry explicit native assurance. Actual client validation passes, including frozen-control and exact-cache checks. A new consumer defers complete independent physical audit during private acceptance and requires it for the selected final state; other private states cannot borrow that proof. Two actual sealed-worker replays match full vectors and all resource fields, reject seven invalid/premature-publication cases, and preserve candidate menus, including nonempty four-candidate sources. Private materialization takes1.34/1.19 s versus7.06/6.97 s for per-candidate independent audit; the selected final independent audit costs5.57 s. These sequential component measurements do not establish full-algorithm speedup. A separate registered live trajectory comparison retains identical selection, search and acceptance.

\subsection{Direct audit avoids another candidate serialization}
Direct auditing consumes the actual native candidate without another serialization. It checks source/request identity, restored timing, terminals, retained edges, untouched networks, replacement nodes, graph edges and complete capacity. Six sealed transactions and six corruption guards pass with zero native calls. BNN checks take21.53/20.79 s versus historical34.73/35.68 s; this is not a paired speedup. Domain/FULLSTA checks remain; forced reservations are unsupported. Automatic trajectory equivalence is still pending.

Reproduction starts from the source and input bindings in \path{multiobjective_completed_seed_transfer_01/completion.json}; its archive preserves all388 observations, including dominated and failed-policy legal candidates. The accompanying projection and feedback receipt distinguish final states, interrupted costs and historical prefixes.

\section{Conclusion}
Current-STA whole transactions expose timing--resource alternatives that strict CPD retention misses. Complete resource-first repair preserves39 fewer SB and two fewer CB at equal CPD for1762 extra wire units; tied-path rejection feedback also enables an exact full-vector improvement on a distinct input. R1-rooted refinement extends the observed low-vertical frontier to33 SB. Matched native controls separate selection and grouping from retained-domain contributions. Further work targets measured validation overhead while preserving actual candidate identity, physical delay, FULL/STA and restoration; broader frozen-policy and complete-runtime measurements remain required.

\begin{thebibliography}{2}
\bibitem{vtr9paper} M. A. Elgammal et al., ``VTR 9: Open-source CAD for fabric and beyond FPGA architecture exploration,'' \emph{ACM Trans. Reconfigurable Technol. Syst.}, 2025, doi:10.1145/3734798.
\bibitem{flatland} C. Hao, A. B. Kahng, B. Pramanik, and I. Youssef, ``Escaping Flatland: A placement flow for enabling 3D FPGAs,'' in \emph{Proc. ACM/IEEE DAC}, 2026, doi:10.1145/3770743.3804424.
\end{thebibliography}
\end{document}
